\chapter{Charge transfer}
\label{ch:ct}

In a charged sample electrons move between atoms: a highly charged ion pulls an electron over the barrier from a
nearby, less charged atom.
The model does this with a classical over-the-barrier criterion (\opt{mcmd-charge-transfer = 1}, the default).

\section{The criterion}

For a pair of atoms of different charge, the less charged one is the donor (charge $Q_D$) and the other the
acceptor ($Q_A$).
The electron goes over the barrier when the pair is closer than
\begin{equation}
  R_{\mathrm{crit}} = \frac{\kc\big(Q_D + 1 + 2\sqrt{(Q_D+1)\,Q_A}\big)}{E_{\mathrm{donor}}},
\end{equation}
with $E_{\mathrm{donor}}$ the binding energy of the donor's outermost electron (\file{energy\_levels\_X.txt}).
Atoms are visited in random order, and each takes the nearest partner that qualifies, within a cutoff of 1.5\,nm
(a cell list keeps the search linear in the number of atoms).

\section{Which shell the electron enters}

Following Boll \emph{et al.}~\cite{boll2022}, the electron enters the acceptor vacancy whose binding energy is closest to
that of the donor orbital it came from.
The binding energy in shell $s$ of configuration $c$ is a difference of total energies,
\begin{equation}
  \varepsilon(c,s) = T(c + e_s) - T(c),
\end{equation}
defined for every shell, also for autoionizing configurations.
This is why \file{total\_energies\_X.txt} is required whenever charge transfer is on.
Only shells up to the ground-state occupation are filled, and a vacancy below an occupied shell (a K hole under a
filled L shell) is skipped: refilling the shell that absorbed the photon would short-circuit the Auger cascade the
rates already describe.

\section{Variants}

\begin{description}
  \item[\opt{mcmd-ct-h-donor-only = 1}] (default): only neutral hydrogen donates, the behaviour of the published
    model.
    Hydrogen barely photoionizes at these photon energies and is stripped early by its neighbours.
  \item[\opt{mcmd-allow-h-ct}] (default 0): allow transfer between two hydrogens.
    This is resonant and moves no net charge; it is off by default and not fully validated.
  \item[\opt{mcmd-charge-transfer-downhill}] (default 0): require the transfer to leave the electron more tightly
    bound, counting the Coulomb energy the pair releases, $(Q_A - Q_D - 1)\kc/R$, by more than $10^{-6}$\,eV.
    On the lysozyme example it changes the number of transfers by an order of magnitude without changing any
    observable: charge transfer saturates, and the extra transfers only reshuffle charge the energetics have already
    fixed.
\end{description}

\section{Scheduling}

The pass over all atoms runs every MD step until nothing has transferred for \opt{mcmd-charge-transfer-idle} steps
(default 2000); it then runs only every \opt{mcmd-charge-transfer-recheck} steps (default 100), and every step again
as soon as a transfer happens.
Because transfer is applied per step rather than at a rate, these two settings change how many transfers happen; they
are not pure performance options.
A negative idle value never goes idle.
